\section{The preregistered holdout study (S1)}\label{sec:s1}

S1 tested the configurations frozen from the main campaign on \SOnePanelN\ scenarios disjoint from it, smoke-tested
before preregistration, with both hosts, two repetitions and a deterministic rule for choosing each host's
configuration. This section gives every preregistered verdict and the figures behind the main paper's S1 numbers.

\paragraph{How the panel was prepared} Before the preregistration was committed, one plain-Sonnet smoke session ran on
each of the \SmokeN\ scenarios (\$\SmokeUSD). It found prompt or grader defects in \SmokeFixed\ scenarios: prompts that
did not state what a check required, and checks that rejected correct answers. These were repaired; genuine agent
mistakes were left as they were; every scenario was then revalidated mechanically (the starting workspace fails, the
reference passes, wrong variants fail). The smoke results enter no hypothesis and are not pooled with S1.

\begin{casecard}{S1 methods in brief}
\textbf{Quality.} Each scripted turn carries automated checks (tests, file patterns, required facts or document
sections); the \term{turn-pass fraction} is the share of a session's turns whose checks passed, and a final-state check
on the hidden tests is reported as additional evidence. The graders were validated before the run: the starting
workspace must fail, the reference solution must pass every turn, and plausible wrong answers must fail.
Non-inferiority means the 95\,\% lower bound of the turn-pass difference is above \NIMargin, that is, no worse than the
specified \NIMarginPts-point margin. These are claims about these automated checks, not about human judgement of, for
example, the quality of a code review or an explanation.

\textbf{Cost.} The primary cost is the tools-normalized list-price cost of each session, plus estimated decider charges
where a decider runs. The normalization reprices the first shared tool-definition prefix from the cache-write to the
cache-read rate; it changed the accepted sessions' recomputed cost by \NormShareMain\,\% in main-v1 and \NormShareSone\,\%
in S1 (data dictionary in the evidence package).

\textbf{Composed policies.} S1 executed pinned-host, pinned-Sonnet and shipped-policy sessions. H1 and H2 evaluate
decide-once policies by selecting the corresponding pinned outcome for each scenario and repetition: the rule or recorded
Jev decision determines which session contributes its cost and quality. These policies therefore reuse executed outcomes
rather than running a separate session for every configuration. Intervals resample whole scenarios. H6 separately
compares live shipped sessions with their matched reconstructions, providing an empirical check of this cost
decomposition (not of every quality property of every reconstructed policy).
\end{casecard}
\begin{table}[htbp]\centering\small
\caption{S1 (holdout-v3) preregistered hypotheses. Cost ratios are geometric means over scenarios; Holm within the two
preregistered families. H3b is descriptive.}\label{tab:s1hyp}
\begin{tabular*}{\textwidth}{@{\extracolsep{\fill}}l r r r l@{}}
\toprule
Hypothesis & cost ratio (95\,\% CI$^{\dagger}$) & turn-pass $\Delta$ (95\,\% CI) & Holm $p$ & verdict \\
\midrule
H1 & 0.581 (0.542--0.625) & $-$0.013 ($-$0.028 to 0.002) & $<$\,0.001 & \ok\ supported \\
H2 & 1.008 (1.000--1.019)$^{\dagger}$ & $-$0.003 ($-$0.008 to 0.000) & $<$\,0.001 & \ok\ supported \\
H3 & 0.999 (0.984--1.017)$^{\dagger}$ & 0.007 ($-$0.004 to 0.022) & $<$\,0.001 & \ok\ supported \\
H3b & 1.004 (0.989--1.019)$^{\dagger}$ & 0.004 ($-$0.002 to 0.010) & -- & descriptive \\
H4 & 0.992 (0.974--1.009) & $-$0.002 ($-$0.013 to 0.009) & 0.187 & \no\ not supported \\
H5 & 1.255 (1.197--1.321) & $-$0.004 ($-$0.024 to 0.017) & $<$\,0.001 & \ok\ supported \\
H6 & 1.014 (0.996--1.032)$^{\dagger}$ & 0.004 ($-$0.011 to 0.019) & $<$\,0.001 & \ok\ supported \\
H7 & 0.616 (0.565--0.672) & $-$0.028 ($-$0.063 to 0.013) & 0.120 & \no\ not supported \\
\bottomrule
\end{tabular*}
\par\smallskip{\footnotesize $^{\dagger}$\,H2, H3, H3b and H6 are equivalence tests; their cost interval is the 90\,\% interval the preregistration uses.}
\end{table}

\ifnum\SOnePending=0
S1 ran \SOneSessions\ sessions on \SOneScen\ scenarios. Table~\ref{tab:s1hyp} lists the verdicts; in plain words:
\begin{itemize}
  \item \textbf{H1 (Fable value), \SHOneVerdict.} The frozen Fable configuration cost \SHOneRatio\X\ plain Fable (95\,\% CI
    \SHOneCI) with turn-pass \SHOneDtp\ (\SHOneDtpCI). The saving is real, but smaller than predicted: the frozen
    prediction from main-v1 was \SPredRatio\ (\SPredCI), and the observed ratio lies above that interval.
  \item \textbf{H2 (decision model against a rule), \SHTwoVerdict.} Letting Jev choose the session's model and letting
    the one-line rule R* choose it were equivalent: \SHTwoRatio\X\ (90\,\% CI \SHTwoCINinety), turn-pass \SHTwoDtp. They
    disagreed on only \SHTwoDisc\ of \SHTwoReps\ scenario-reps. The preregistered decision follows: ship the simpler rule
    R*, with no external call and no consent needed; Jev becomes an opt-in decider.
  \item \textbf{H3 (bundle overhead), \SHThreeVerdict.} The shipped bundle on Opus cost \SHThreeRatio\X\ plain Opus (90\,\%
    CI \SHThreeCINinety): cost-equivalent within the preregistered $\pm$5\,\% margin. On Fable, the bundle with the host pinned cost \SHThreeBRatio\X\ (descriptive).
  \item \textbf{H4 (Opus at medium effort), \SHFourVerdict.} Opus at medium cost \SHFourRatio\X\ its default (95\,\% CI
    \SHFourCI): medium effort did not demonstrate a saving, unlike Sonnet (\EcRatio\X) and Fable (\EfRatio\X).
  \item \textbf{H5 (routing on Opus), \SHFiveVerdict.} Routing to Sonnet still lost on Opus: \SHFiveRatio\X\ (\SHFiveCI),
    on fresh scenarios and without the effort-switch artefact.
  \item \textbf{H6 (live equals predicted), \SHSixVerdict.} Live decide-once sessions cost what the decomposition predicted
    (\SHSixRatio\X, 90\,\% CI \SHSixCINinety).
  \item \textbf{H7 (review and explain on Sonnet), \SHSevenVerdict.} On the \SHSevenQualN\ review and explain scenarios
    (\SHSevenN\ for cost, after the gate-failed scenario is excluded), the turn-pass lower bound for Sonnet against the host was \SHSevenDtpLow, below the \NIMargin\ margin: non-inferiority was not
    established (the interval also includes zero, so a loss is not shown either). The preregistered
    rule therefore keeps these task types on the host (\code{keep\_on\_host: [review, explain]}).
\end{itemize}

\paragraph{What the freeze rule chose} On Fable the rule picked the configuration (\SFreezeFableRatio\X)
\begin{center}\small\code{\SFreezeFable}\end{center}
not the preregistered
\begin{center}\small\code{\SFreezeFablePrereg}\end{center} That makes it \emph{selected on holdout}: it must replicate in a follow-up study before it
ships. It also conflicts with H7, which keeps review and explain on the host. What ships for Fable until S2 replicates is
therefore the preregistered configuration plus \code{keep\_on\_host: [review, explain]}, applied through a keyword proxy
(main paper, Section~\ref{M-sec:recs}). On Opus no candidate qualified, so
Opus keeps the shipped defaults: no routing, default effort.

\begin{keyidea}
The decision model's start-of-session choice was matched by a one-line rule on these hosts and prices. Where the money
is, routing on Fable and never on Opus, price arithmetic and a rule decide it. The value of a system-1 model is not in
this decision.
\end{keyidea}

\begin{processnote}{Process note: how the S1 run went}
\raggedright The study spent \$\SSpend. \SGateFail\ Fable sessions of one scenario
(a code-review scenario) were served partly by Opus, most likely the provider's fallback (an inference); the preregistered rule
excluded them from cost, so Fable cost contrasts use \SHOneN\ scenarios. The stop rule for gate failures was not followed
to the letter, because the gate is evaluated when rows are extracted, not live; this is a disclosed deviation. A Forge
outage made \SInfraAttempts\ of \SAttempts\ wave attempts (\SInfraPct\,\%) fail on infrastructure; every failed wave was
rerun whole. The budget cap was raised from \$\SBudgetOrig\ to \$\SBudgetFinal\ during the run, though the final spend
stayed under the original cap, and concurrency changed several times. The A/A copies differed in turn-pass by
\SNoiseOpusDtp\ on Opus (\SNoiseN\ pairs), so small quality differences cannot be interpreted. Rerunning the analysis
without the flagged scenarios changed no verdict and no freeze choice.
\end{processnote}

\subsubsection{Every hypothesis on one page}

\Cref{fig:s1forest} puts all eight preregistered comparisons side by side, cost on the left and quality on the right, so the reader can see at once which results move money and which are guards.

\begin{figure}[htbp]
\centering
\pgfplotslegendfromname{sforestlegend}\par\smallskip
% S1: every preregistered hypothesis, cost panel (log) and quality panel.
\begin{tikzpicture}
\begin{axis}[benchmark, name=cost, xmode=log, log ticks with fixed point, xtick={0.6,0.8,1,1.2}, xmin=0.5, xmax=1.4,
  scale only axis, width=5.0cm, height=6.2cm, xmajorgrids, ymin=-0.6, ymax=7.95, ytick={0,...,7},
  yticklabels from table={generated/data/s1-forest.dat}{label}, y tick label style={font=\footnotesize},
  xlabel={Cost ratio (log)}, title={Cost}, title style={font=\small\bfseries},
  legend to name=sforestlegend, legend columns=3, legend cell align=left,
  legend style={/tikz/every even column/.append style={column sep=10pt}}]
\draw[black!70, line width=0.8pt] (axis cs:1,-0.6) -- (axis cs:1,7.6);
\draw[okorange, dashed, line width=0.8pt] (axis cs:\SThreshOne,-0.6) -- (axis cs:\SThreshOne,7.6);
\addplot[nodes near coords, point meta=explicit symbolic, nodes near coords style={font=\scriptsize, fill=white, fill opacity=0.85, text opacity=1, inner sep=1pt, anchor=south, yshift=2.5pt, text=black!80}, only marks, mark=*, mark size=2.4pt, color=okblue, error bars/.cd, x dir=both, x explicit,
  error bar style={line width=1pt}] table[meta=v, x=r, y=y, x error minus=rm, x error plus=rp]{generated/data/s1-forest.dat};
\addlegendentry{estimate, 95\,\% CI}
\addlegendimage{black!70, line width=0.8pt}\addlegendentry{no difference}
\addlegendimage{okorange, dashed, line width=0.8pt}\addlegendentry{preregistered threshold}
\end{axis}
\begin{axis}[benchmark, at={(cost.south east)}, xshift=0.5cm, scale only axis, width=4.0cm, height=6.2cm,
  xmin=-0.09, xmax=0.06, xtick={-0.05,0,0.05}, scaled x ticks=false, xticklabel style={/pgf/number format/fixed}, xmajorgrids,
  ymin=-0.6, ymax=7.95, ytick={0,...,7}, yticklabels={}, xlabel={Turn-pass difference}, title={Quality},
  title style={font=\small\bfseries}]
\draw[black!70, line width=0.8pt] (axis cs:0,-0.6) -- (axis cs:0,7.6);
\draw[okorange, dashed, line width=0.8pt] (axis cs:-0.05,-0.6) -- (axis cs:-0.05,7.6);
\addplot[only marks, mark=square*, mark size=2.2pt, color=okgreen, error bars/.cd, x dir=both, x explicit,
  error bar style={line width=1pt}] table[x=d, y=y, x error minus=dm, x error plus=dp]{generated/data/s1-forest.dat};
\end{axis}
\end{tikzpicture}

\caption{Two hypotheses carry money (H1, H5); the rest are guards. S1 hypotheses H1--H7: cost ratio (left, log scale) and turn-pass difference (right), each with its 95\,\%
interval. H3b is descriptive. Thresholds: the H1 cost bound (\SThreshOne) and the quality margin (\NIMargin).}
\label{fig:s1forest}
\howtoread{Read each row across both panels. A row helps a configuration when its cost dot is left of 1 and its
quality whisker stays right of the dashed line. H1 does both. H5 is far right of 1: routing on Opus costs more. H4
straddles 1: medium effort did nothing measurable on Opus. H7's quality whisker crosses the dashed line, which is why
review and explain stay on the host.}
\end{figure}

\begin{keyidea}
Two of the eight rows carry money: H1 (Fable routing saves about two fifths) and H5 (Opus routing loses about a
quarter). The rest are guards: the bundle is cost-equivalent within the \(\pm\)5\,\% margin (H3), the model adds nothing over the rule (H2), live sessions behave as
predicted (H6), and review and explain need protection (H7).
\end{keyidea}

\subsubsection{How to read a paired cost ratio}\label{sec:worked}

\begin{casecard}{Worked example: one scenario, two reps, one ratio}
Take one explain scenario from the holdout study, the scenario whose ratio sits closest to the overall result. Under the frozen Fable
configuration it was \WxArm. Across its \WxReps\ repetitions the routed sessions cost \$\WxRouted\ on average and the
plain Fable sessions \$\WxPlain. Its paired cost ratio is therefore $\WxRouted / \WxPlain = \WxRatio$; on the log scale
that is $\ln \WxRatio$ = \WxLog. Doing the same for each of the \WxNScen\ scenarios, averaging the logs and
exponentiating gives the \term{geometric mean}, \WxGM: the H1 headline. The interval comes from resampling whole
scenarios \SResamples\ times and repeating the calculation.
\end{casecard}

Why a geometric mean? A ratio of 0.5 and a ratio of 2 should cancel; their arithmetic mean (1.25) says ``more
expensive'', their geometric mean (1) says ``no change''. Averaging logs treats halving and doubling symmetrically.

\subsubsection{The variance underneath the headline}

A geometric mean hides how much individual tasks vary. \Cref{fig:perscen} shows every scenario's paired ratio, sorted, for the two results that carry money.

\begin{figure}[htbp]
\centering
\pgfplotslegendfromname{pslegend}\par\smallskip
% S1: per-scenario paired cost ratios, sorted (log y).
\begin{tikzpicture}
\begin{axis}[benchmark, name=f, ymode=log, log ticks with fixed point, ytick={0.4,0.6,0.8,1,1.5,2}, ymin=0.33, ymax=2.4,
  scale only axis, width=5.6cm, height=5cm, ymajorgrids, xmin=0, xmax=61, xlabel={Scenario (sorted)},
  ylabel={Cost ratio (log)}, title={Fable: frozen config / plain}, title style={font=\small\bfseries},
  legend to name=pslegend, legend columns=2, legend cell align=left,
  legend style={/tikz/every even column/.append style={column sep=10pt}}]
\draw[black!70, line width=0.8pt] (axis cs:0,1) -- (axis cs:61,1);
\addplot[only marks, mark=*, mark size=1.6pt, color=okblue] table[x=x, y=r]{generated/data/perscen-fable.dat};
\addlegendentry{one scenario (mean of its reps)}
\addplot[only marks, mark=triangle*, mark size=1.9pt, color=okorange] coordinates {(-10,1)};
\addlegendentry{Opus panel}
\end{axis}
\begin{axis}[benchmark, at={(f.south east)}, xshift=1.2cm, ymode=log, log ticks with fixed point,
  ytick={0.4,0.6,0.8,1,1.5,2}, ymin=0.33, ymax=2.4, scale only axis, width=5.6cm, height=5cm, ymajorgrids,
  xmin=0, xmax=61, xlabel={Scenario (sorted)}, title={Opus: Sonnet / plain}, title style={font=\small\bfseries}]
\draw[black!70, line width=0.8pt] (axis cs:0,1) -- (axis cs:61,1);
\addplot[only marks, mark=triangle*, mark size=1.9pt, color=okorange] table[x=x, y=r]{generated/data/perscen-opus.dat};
\end{axis}
\end{tikzpicture}

\caption{The direction is consistent across scenarios; the size varies widely. Per-scenario paired cost ratios in S1, sorted. Left: the frozen Fable configuration against plain Fable
(\PsFableN\ scenarios). Right: Sonnet decided once against plain Opus (\PsOpusN\ scenarios). Log scale.}
\label{fig:perscen}
\howtoread{Each dot is one scenario. Below the line saves money. On Fable, \PsFableAbove\ of \PsFableN\ scenarios cost
more than plain Fable, and the ratios run from \PsFableMin\ to \PsFableMax. On Opus, \PsOpusAbove\ of \PsOpusN\ cost
more, from \PsOpusMin\ to \PsOpusMax. The direction is consistent; the size varies a lot from task to task.}
\end{figure}

\subsubsection{Where the risk sits: task types and families}

The quality cost of routing is not spread evenly. \Cref{fig:subgroups} splits the Fable result by task type and by scenario family (exploratory).

\begin{figure}[htbp]
\centering
\pgfplotslegendfromname{sublegend}\par\smallskip
% S1, Fable: the frozen configuration against plain Fable, per task type and per family.
\begin{tikzpicture}
\begin{axis}[benchmark, name=cost, xmode=log, log ticks with fixed point, xtick={0.4,0.6,0.8,1}, xmin=0.35, xmax=1.1,
  scale only axis, width=5.0cm, height=6.6cm, xmajorgrids, ymin=-0.6, ymax=10.95, ytick={\SubTicks},
  yticklabels/.expanded={\SubLabels}, y tick label style={font=\footnotesize},
  xlabel={Cost ratio (log)}, title={Cost}, title style={font=\small\bfseries},
  legend to name=sublegend, legend columns=2, legend cell align=left,
  legend style={/tikz/every even column/.append style={column sep=10pt}}]
\draw[black!70, line width=0.8pt] (axis cs:1,-0.6) -- (axis cs:1,10.6);
\addplot[nodes near coords, point meta=explicit symbolic, nodes near coords style={font=\scriptsize, fill=white, fill opacity=0.85, text opacity=1, inner sep=1pt, anchor=south, yshift=2.5pt, text=black!80}, only marks, mark=*, mark size=2.4pt, color=okblue, error bars/.cd, x dir=both, x explicit,
  error bar style={line width=1pt}] table[meta=v, x=r, y=y, x error minus=rm, x error plus=rp]{generated/data/s1-subgroups.dat};
\addlegendentry{estimate, 95\,\% CI}
\addlegendimage{okorange, dashed, line width=0.8pt}\addlegendentry{quality margin}
\end{axis}
\begin{axis}[benchmark, at={(cost.south east)}, xshift=0.5cm, scale only axis, width=4.0cm, height=6.6cm,
  xmin=-0.12, xmax=0.07, xtick={-0.1,-0.05,0,0.05}, scaled x ticks=false, xticklabel style={/pgf/number format/fixed}, xmajorgrids,
  ymin=-0.6, ymax=10.95, ytick={\SubTicks}, yticklabels={}, xlabel={Turn-pass difference}, title={Quality},
  title style={font=\small\bfseries}]
\draw[black!70, line width=0.8pt] (axis cs:0,-0.6) -- (axis cs:0,10.6);
\draw[okorange, dashed, line width=0.8pt] (axis cs:-0.05,-0.6) -- (axis cs:-0.05,10.6);
\addplot[only marks, mark=square*, mark size=2.2pt, color=okgreen, error bars/.cd, x dir=both, x explicit,
  error bar style={line width=1pt}] table[x=d, y=y, x error minus=dm, x error plus=dp]{generated/data/s1-subgroups.dat};
\end{axis}
\end{tikzpicture}

\caption{Every subgroup saves money; quality intervals for several subgroups are too wide to rule out a loss. Lower bounds fall below the $-$0.05 margin for the explain (\SubExplainTaskLow), mixed (\SubMixedTaskLow) and docs (\SubDocsTaskLow) tasks and the mixed (\SubMixedFamilyLow) and knowledge (\SubKnowledgeFamilyLow) families; review's lower bound is \SubReviewTaskLow. Subgroups hold 9--20 scenarios; treat as hypothesis-generating. The frozen Fable configuration against plain Fable, split by task type and by scenario family (number of
cost-valid scenarios in brackets; the quality panel is unfiltered and uses every scenario of the subgroup), with 95\,\%
intervals. Exploratory: subgroups were not preregistered except H7.}
\label{fig:subgroups}
\howtoread{Every subgroup saves money. The quality panel is where they differ: explain (\SubExplainTaskDtp; lower
bound \SubExplainTaskLow) and mixed work cross the margin, review touches it (\SubReviewTaskLow), while bugfix and
feature sit near zero. Small subgroups give wide whiskers; treat single rows with caution.}
\end{figure}

\subsubsection{Which configuration the rule picked}

The freeze rule considered every configuration the preregistration listed and chose one per host. \Cref{fig:fz} shows them all, with the choice and the preregistered configuration marked; Table~\ref{tab:fzkey} lists them.

\begin{figure}[htbp]
\centering
\pgfplotslegendfromname{fzlegendFable}\par\smallskip
\fzpanel{Fable}{Fable}{fable}\par\medskip
\fzpanel{Opus}{Opus}{opus}
\caption{On Fable many configurations qualify; on Opus none does. All \FzFableN\ Fable and \FzOpusN\ Opus configurations the freeze rule considered: cost ratio against the
host's plain default and turn-pass difference (point estimates). Keys group configurations that land on the same
point; Table~\ref{tab:fzkey} lists them.}
\label{fig:fz}
\howtoread{Left and up is better. Filled dots passed the preregistered candidate test (cost upper bound below 1,
quality lower bound above the margin, final-state point no worse than the specified limit). The star is the rule's choice; the
diamond is the preregistered configuration. On Fable \FzFableCand\ configurations qualified; the cheapest cluster
drops the scope gate. On Opus none qualified, so the shipped default stays.}
\end{figure}

\subsubsection{The decision model against the rule}\label{sec:bound}

Jev and the rule R* chose differently on only \DcN\ of \DcUnits\ Fable scenario-reps (\cref{tab:discordant}).

\begin{table}[htbp]
\centering\small
\caption{The \DcN\ S1 scenario-reps on which the decision model (Jev) and the rule R* chose different session
models, with what each choice cost and how it scored.}
\label{tab:discordant}
\begin{tabular*}{\textwidth}{@{\extracolsep{\fill}}l r l l r r r r@{}}
\toprule
 & & \multicolumn{2}{c}{model chosen by} & \multicolumn{2}{c}{session cost} & \multicolumn{2}{c}{turn-pass} \\
\cmidrule(lr){3-4}\cmidrule(lr){5-6}\cmidrule(l){7-8}
Scenario & rep & Jev & rule R* & Jev & R* & Jev & R* \\
\midrule
go-bowling & 1 & Fable, medium & Sonnet (Pc) & \$6.53 & \$4.06 & 0.857 & 1.000 \\
go-bowling & 2 & Fable, medium & Sonnet (Pc) & \$5.66 & \$4.77 & 0.929 & 0.929 \\
peewee-filterfk & 1 & Fable, medium & Sonnet (Pc) & \$2.15 & \$1.94 & 0.889 & 1.000 \\
peewee-filterfk & 2 & Fable, medium & Sonnet (Pc) & \$2.34 & \$1.87 & 0.889 & 1.000 \\
\bottomrule
\end{tabular*}

\end{table}

\begin{casecard}{Worked example: why the two deciders cannot differ much}
Two deciders that agree on a scenario-rep produce the same session there, so their difference comes only from the
scenario-reps where they disagree. The share of those is $d = \DcN/\DcUnits = \DcD$. On them the choices differed in
cost by \$\DcMeanAbsDiff\ per session on average. So the most the deciders can differ, averaged over all sessions, is
$d \times \$\DcMeanAbsDiff = \$\DcBound$ per session, about \DcBoundPct\,\% of a typical session (\$\DcMeanCost).
That bound, not a significance test, is what makes ``equivalent'' believable: with so few disagreements, no amount of
skill on those few can move the total.
\end{casecard}

\subsubsection{Effort, across three hosts}

Three separate studies measured medium against default reasoning effort, one per host model. \Cref{fig:effort} puts them on one chart.

\begin{figure}[htbp]
\centering
\pgfplotslegendfromname{efflegend}\par\smallskip
% Medium vs default reasoning effort, per host, three studies.
\begin{tikzpicture}
\begin{axis}[benchmark, name=cost, scale only axis, width=5.2cm, height=2.8cm, xmin=0.75, xmax=1.05,
  xtick={0.8,0.9,1}, xmajorgrids, ymin=-0.6, ymax=2.95, ytick={0,1,2},
  yticklabels from table={generated/data/effort-hosts.dat}{label}, y tick label style={font=\footnotesize},
  xlabel={Cost ratio, medium / default}, title={Cost}, title style={font=\small\bfseries},
  legend to name=efflegend, legend columns=2, legend cell align=left,
  legend style={/tikz/every even column/.append style={column sep=10pt}}]
\draw[black!70, line width=0.8pt] (axis cs:1,-0.6) -- (axis cs:1,2.6);
\addplot[nodes near coords, point meta=explicit symbolic, nodes near coords style={font=\scriptsize, fill=white, fill opacity=0.85, text opacity=1, inner sep=1pt, anchor=south, yshift=2.5pt, text=black!80}, only marks, mark=*, mark size=2.4pt, color=okblue, error bars/.cd, x dir=both, x explicit,
  error bar style={line width=1pt}] table[meta=v, x=r, y=y, x error minus=rm, x error plus=rp]{generated/data/effort-hosts.dat};
\addlegendentry{cost, 95\,\% CI}
\addplot[only marks, mark=square*, mark size=2.2pt, color=okgreen] coordinates {(5,5)};
\addlegendentry{turn-pass, 95\,\% CI}
\end{axis}
\begin{axis}[benchmark, at={(cost.south east)}, xshift=0.5cm, scale only axis, width=4.0cm, height=2.8cm,
  xmin=-0.06, xmax=0.09, xtick={-0.05,0,0.05}, scaled x ticks=false, xticklabel style={/pgf/number format/fixed}, xmajorgrids,
  ymin=-0.6, ymax=2.95, ytick={0,1,2}, yticklabels={}, xlabel={Turn-pass difference}, title={Quality},
  title style={font=\small\bfseries}]
\draw[black!70, line width=0.8pt] (axis cs:0,-0.6) -- (axis cs:0,2.6);
\draw[okorange, dashed, line width=0.8pt] (axis cs:-0.05,-0.6) -- (axis cs:-0.05,2.6);
\addplot[only marks, mark=square*, mark size=2.2pt, color=okgreen, error bars/.cd, x dir=both, x explicit,
  error bar style={line width=1pt}] table[x=d, y=y, x error minus=dm, x error plus=dp]{generated/data/effort-hosts.dat};
\end{axis}
\end{tikzpicture}

\caption{Medium effort saved money on Sonnet and Fable but did not demonstrate a saving on Opus. Medium against default reasoning effort, measured in three studies: Sonnet (follow-up 1), Fable (follow-up
2) and Opus (S1, H4). 95\,\% intervals.}
\label{fig:effort}
\howtoread{Sonnet and Fable sit well left of 1: medium effort is cheaper with no detectable quality loss. Opus sits
on 1. The same setting saves money on two models and nothing on the third, so effort must be set per host.}
\end{figure}

\subsubsection{Where the money goes}

Why does the same switch save money on one host and cost money on the other? \Cref{fig:composition} splits the average S1 session's cost by token class.

\begin{figure}[htbp]
\centering
\pgfplotslegendfromname{complegend}\par\smallskip
% S1: mean dollars per session by token class.
\begin{tikzpicture}
\begin{axis}[benchmark, xbar stacked, bar width=10pt, scale only axis, width=10cm, height=4.2cm, xmin=0, xmajorgrids,
  ymin=-0.6, ymax=3.6, ytick={0,1,2,3}, yticklabels from table={generated/data/s1-composition.dat}{label},
  y tick label style={font=\small}, xlabel={Mean cost per session (US dollars)},
  title={Cost per session, by token class}, title style={font=\small\bfseries},
  legend to name=complegend, legend columns=4, legend cell align=left,
  legend style={/tikz/every even column/.append style={column sep=8pt}}]
\addplot[fill=okgray!60, draw=black!60] table[x=inp, y=y]{generated/data/s1-composition.dat};
\addlegendentry{uncached input}
\addplot[fill=okblue!55, draw=black!60] table[x=rd, y=y]{generated/data/s1-composition.dat};
\addlegendentry{cache read}
\addplot[fill=okorange!80, draw=black!60, postaction={pattern=north east lines, pattern color=white}] table[x=wr, y=y]{generated/data/s1-composition.dat};
\addlegendentry{cache write}
\addplot[fill=okgreen!60, draw=black!60, postaction={pattern=dots, pattern color=white}] table[x=out, y=y]{generated/data/s1-composition.dat};
\addlegendentry{output}
\node[anchor=west, font=\scriptsize, text=black!80] at (axis cs:1.9122,0) {1.91};
\node[anchor=west, font=\scriptsize, text=black!80] at (axis cs:2.3392,1) {2.34};
\node[anchor=west, font=\scriptsize, text=black!80] at (axis cs:3.0271,2) {3.03};
\node[anchor=west, font=\scriptsize, text=black!80] at (axis cs:4.5530,3) {4.55};

\end{axis}
\end{tikzpicture}

\caption{Plain Fable spends most of its money writing to the cache. Mean cost per S1 session split by token class, at the study's price table (exploratory).}
\label{fig:composition}
\howtoread{Each bar is the average session. Plain Fable spends most on cache writes (\CcPlainFableWriteShare\,\% of
\$\CcPlainFableTotal): writing context into Fable's cache is expensive. The Sonnet session (\$\CcSonnetDecidedOTotal)
spends more on cache reads (\CcSonnetDecidedOReadShare\,\%), which are cheap per token but many. Plain Opus
(\$\CcPlainOpusTotal) is already cheaper than the Sonnet session, which is the whole H5 story in one bar.}
\end{figure}


\subsubsection{Every configuration the freeze rule considered}


\begingroup\small
\FloatBarrier
\begin{xltabular}{\textwidth}{l l >{\raggedright\arraybackslash}X S[table-format=1.3] S[table-format=-1.3] c}
\caption{Every configuration the S1 freeze rule considered (keys as in \cref{fig:fz}). Cost ratio and turn-pass difference against the host's plain default.}\label{tab:fzkey}\\
\toprule Key & Host & Configuration & {cost ratio} & {turn-pass $\Delta$} & candidate \\ \midrule \endfirsthead
\toprule Key & Host & Configuration & {cost ratio} & {turn-pass $\Delta$} & candidate \\ \midrule \endhead
F1 & Fable & always\_route | strong default | scope off | keep none & 0.528 & -0.017 & \ok \\
F1 & Fable & always\_route | strong medium | scope off | keep none & 0.528 & -0.017 & \ok \\
F1 & Fable & rstar | strong default | scope off | keep none & 0.528 & -0.017 & \ok \\
F1 & Fable & rstar | strong medium | scope off | keep none & 0.528 & -0.017 & \ok \\
F2 & Fable & always\_route | strong medium | scope 300 | keep none & 0.581 & -0.013 & \ok \\
F2 & Fable & rstar | strong medium | scope 300 | keep none & 0.581 & -0.013 & \ok \\
F2 & Fable & jev | strong medium | scope 300 | keep none & 0.586 & -0.016 & \ok \\
F3 & Fable & always\_route | strong default | scope 300 | keep none & 0.608 & -0.012 & \ok \\
F3 & Fable & rstar | strong default | scope 300 | keep none & 0.608 & -0.012 & \ok \\
F3 & Fable & jev | strong default | scope 300 | keep none & 0.618 & -0.015 & \ok \\
F4 & Fable & always\_route | strong medium | scope off | keep review+explain & 0.618 & -0.007 & \ok \\
F4 & Fable & rstar | strong medium | scope off | keep review+explain & 0.618 & -0.007 & \ok \\
F5 & Fable & always\_route | strong medium | scope 300 | keep review+explain & 0.650 & -0.006 & \ok \\
F5 & Fable & rstar | strong medium | scope 300 | keep review+explain & 0.650 & -0.006 & \ok \\
F5 & Fable & jev | strong medium | scope 300 | keep review+explain & 0.655 & -0.009 & \ok \\
F6 & Fable & always\_route | strong default | scope off | keep review+explain & 0.650 & -0.001 & \ok \\
F6 & Fable & rstar | strong default | scope off | keep review+explain & 0.650 & -0.001 & \ok \\
F7 & Fable & always\_route | strong default | scope 300 | keep review+explain & 0.706 & -0.000 & \ok \\
F7 & Fable & rstar | strong default | scope 300 | keep review+explain & 0.706 & -0.000 & \ok \\
F7 & Fable & jev | strong default | scope 300 | keep review+explain & 0.717 & -0.003 & \ok \\
F8 & Fable & plain-medium & 0.851 & 0.005 & \ok \\
F9 & Fable & bundle-unrouted medium & 0.853 & -0.005 & \ok \\
F10 & Fable & bundle-unrouted default & 1.004 & 0.004 &  \\
O1 & Opus & plain-medium & 0.992 & -0.002 &  \\
O2 & Opus & bundle-unrouted default & 0.999 & 0.007 &  \\
O2 & Opus & bundle-unrouted medium & 1.009 & 0.003 &  \\
O3 & Opus & jev | strong default | scope 300 | keep review+explain & 1.080 & 0.005 &  \\
O4 & Opus & jev | strong medium | scope 300 | keep review+explain & 1.091 & 0.000 &  \\
O5 & Opus & always\_route | strong default | scope 300 | keep review+explain & 1.097 & 0.005 &  \\
O5 & Opus & rstar | strong default | scope 300 | keep review+explain & 1.097 & 0.005 &  \\
O6 & Opus & always\_route | strong medium | scope 300 | keep review+explain & 1.109 & 0.004 &  \\
O6 & Opus & rstar | strong medium | scope 300 | keep review+explain & 1.109 & 0.004 &  \\
O7 & Opus & jev | strong default | scope 300 | keep none & 1.136 & -0.002 &  \\
O7 & Opus & jev | strong medium | scope 300 | keep none & 1.144 & -0.006 &  \\
O8 & Opus & always\_route | strong default | scope 300 | keep none & 1.153 & -0.002 &  \\
O8 & Opus & rstar | strong default | scope 300 | keep none & 1.153 & -0.002 &  \\
O8 & Opus & always\_route | strong medium | scope 300 | keep none & 1.163 & -0.002 &  \\
O8 & Opus & rstar | strong medium | scope 300 | keep none & 1.163 & -0.002 &  \\
O9 & Opus & always\_route | strong default | scope off | keep review+explain & 1.182 & 0.005 &  \\
O9 & Opus & rstar | strong default | scope off | keep review+explain & 1.182 & 0.005 &  \\
O9 & Opus & always\_route | strong medium | scope off | keep review+explain & 1.187 & 0.005 &  \\
O9 & Opus & rstar | strong medium | scope off | keep review+explain & 1.187 & 0.005 &  \\
O10 & Opus & always\_route | strong default | scope off | keep none & 1.255 & -0.004 &  \\
O10 & Opus & always\_route | strong medium | scope off | keep none & 1.255 & -0.004 &  \\
O10 & Opus & rstar | strong default | scope off | keep none & 1.255 & -0.004 &  \\
O10 & Opus & rstar | strong medium | scope off | keep none & 1.255 & -0.004 &  \\
\bottomrule
\end{xltabular}

\FloatBarrier
\endgroup


